Recursive synthetic training, where each model generation fine-tunes on its predecessor's outputs, causes model collapse: a progressive loss of distributional diversity and downstream capability.
Detecting collapse currently requires costly text generation or benchmark evaluation at every generation.
Spectral analysis of internal representations offers a cheaper alternative, but the spectral dynamics of recursive collapse lack systematic characterization, and no prior work connects spectral indicators to closed-loop prevention.
We track spectral properties of the embedding Gram matrix across 10 recursive generations of Pythia-410M, with pilot validation on Pythia-1B.
Our analysis reveals a reproducible three-phase pattern: spectral decompression (generations 0--2), where effective rank and spectral log-volume (SLV) increase even as downstream performance degrades; collapse onset (generations 2--4), where spectral indicators reverse; and accelerating collapse (generations 4--10), where spectral contraction proceeds at 2.43$\times$ the onset rate.
Among the indicators we examine, SLV is particularly suited for automated monitoring: it requires only a forward pass on held-out data, captures distributed volume changes across all spectral directions, and exhibits post-decompression monotonic decline amenable to threshold-based control.
We further present a proof-of-concept proportional controller that uses SLV to guide real-synthetic data mixing, and identify a signal attenuation paradox where successful intervention suppresses the spectral signal the controller relies on.